\subsection{Independent four-block closure classes}
\label{sec:peps_dependent_closure_dimension}

Retain the four blocks, eight separately represented bonds, and actual cut
spaces of Definition~\ref{def:peps_full_four_cut_space}. Write
$p^{g,h}$ for the standard seam labels and $A^0$ for the four incident
averaging tensors. Let $K$ be the number of simultaneous conjugacy classes
of commuting pairs in $G$. The independence argument of
\cite[Theorem~5.9]{Schuch2010PEPS} gives an exact dimension for this
four-block intersection. The identification with a microscopic parent
kernel is a separate geometric assertion.
% Scope: this fragment proves closure independence and four-cut dimension.
% It must not replace the unrestricted microscopic source-labelled theorem.
% Source: Papers/1001.3807/paper_v3.tex:1560--1621.

\begin{theorem}[Synchronization of closure gauges]%
    \label{thm:peps_dependent_closure_gauge}
    \lean{TNLean.PEPS.DependentTorus.closureLabels_conjugate,
        TNLean.PEPS.DependentTorus.nonseamCompatible_of_closureLabels_eq,
        TNLean.PEPS.DependentTorus.closureLabels_eq_iff}
    \leanok
    \uses{def:peps_labelled_closure}
    For any $x,g,h\in G$, one has
    $p^{xgx^{-1},xhx^{-1}}_e=xp^{g,h}_ex^{-1}$ on every bond.
    Given pairs $(a,b),(g,h)$ and vertex labels $q$, the identities
    \begin{align}
        p^{a,b}_e=q_{\operatorname{head}(e)}p^{g,h}_e
        q_{\operatorname{tail}(e)}^{-1}
        \label{eq:peps_dependent_closure_gauge}
    \end{align}
    hold on all eight bonds if and only if $q$ is constant, with value
    $x=q_{(0,0)}$, and $a=xgx^{-1}$ and $b=xhx^{-1}$.
    Commutation is unnecessary here.
\end{theorem}
\begin{proof}\leanok
    \uses{thm:peps_torus_nonseam_connectivity}
    Both standard closures equal $1$ off their seams. On those bonds the
    asserted identities give equality of neighboring vertex labels.
    Nonseam connectivity implies $q_v=q_{(0,0)}$ at all four vertices.
    Reading the vertical bond based at $(0,-1)$ and the horizontal bond
    based at $(-1,0)$ gives $a=xgx^{-1}$ and $b=xhx^{-1}$.
    Conversely, these equalities give the seam identities, while constant
    conjugation fixes the identity labels on the other bonds.
\end{proof}

\begin{theorem}[Exact trace-dual closure coefficients]%
    \label{thm:peps_dependent_closure_extraction}
    \lean{TNLean.PEPS.DependentTorus.bondCoefficientExtraction_closure,
        TNLean.PEPS.DependentTorus.bondCoefficientExtraction_closure_eq_zero_of_class_ne,
        TNLean.PEPS.DependentTorus.bondCoefficientExtraction_closure_self,
        TNLean.PEPS.DependentTorus.bondCoefficientExtraction_closure_self_ne_zero}
    \leanok
    \uses{def:peps_dependent_bond_coefficient, def:peps_full_four_cut_canonical,
        def:peps_pair_conjugacy, def:peps_semiregular}
    Suppose $G$ is finite and every $U_e$ is semi-regular. Let
    $E_{a,b}=c_{p^{a,b}}$ be the product of the actual eight bond
    trace-dual functionals. Then
    \begin{align}
        E_{a,b}\bigl(\zeta_{A^0}(g,h)\bigr)
         & = |G|^{-4}\sum_{x\in G}
        \mathbf 1_{a=xgx^{-1},\,b=xhx^{-1}},
        \label{eq:peps_dependent_closure_extract} \\
        E_{g,h}\bigl(\zeta_{A^0}(g,h)\bigr)
         & = |G|^{-4}|C_G(g,h)|\ne0.
        \label{eq:peps_dependent_closure_diagonal}
    \end{align}
    Here $C_G(g,h)=\{x\in G:xg=gx,\ xh=hx\}$ is the common centralizer.
    In particular, the coefficient is zero for different simultaneous
    conjugacy classes. These formulas hold even for noncommuting pairs.
\end{theorem}
\begin{proof}\leanok
    \uses{thm:peps_dependent_closure_gauge, thm:peps_dependent_canonical_bond_sum}
    Expanding the four normalized local averages and applying the eight
    trace-dual functionals gives
    \begin{align}
        E_{a,b}\bigl(\zeta_{A^0}(g,h)\bigr)
        = |G|^{-4}\sum_{q\in G^\Lambda}
        \prod_e\mathbf 1_{p^{a,b}_e
                =q_{\operatorname{head}(e)}p^{g,h}_e
                q_{\operatorname{tail}(e)}^{-1}}.
        \label{eq:peps_dependent_closure_gauge_sum}
    \end{align}
    Theorem~\ref{thm:peps_dependent_closure_gauge} identifies the surviving
    labels with constant functions $q_v=x$, proving
    (\ref{eq:peps_dependent_closure_extract}). On the diagonal the surviving
    $x$ are precisely the common centralizer. It contains the identity,
    so (\ref{eq:peps_dependent_closure_diagonal}) is nonzero.
\end{proof}

\begin{theorem}[Conjugacy invariance of the actual closures]%
    \label{thm:peps_dependent_closure_invariant}
    \lean{TNLean.PEPS.DependentTorus.closure_conjugate,
        TNLean.PEPS.DependentTorus.closure_eq_of_pairConjugacyClass_eq}
    \leanok
    \uses{def:peps_full_four_cut_space, def:peps_pair_conjugacy}
    For a finite group and locally invariant tensors $A_v$, one has
    $\zeta_A(xgx^{-1},xhx^{-1})=\zeta_A(g,h)$ for all $x,g,h$.
    Thus closures depend only on their simultaneous conjugacy class.
    Neither semi-regularity nor injectivity is needed for this assertion.
\end{theorem}
\begin{proof}\leanok
    \uses{thm:peps_dependent_closure_gauge, thm:peps_dependent_vertex_gauge,
        thm:peps_full_four_cut_transport}
    A constant vertex gauge transforms the two canonical seam labels by
    simultaneous conjugation and leaves the canonical network unchanged.
    Writing $T_v$ for the original site maps, local invariance gives
    $T_vP_v=T_v$. Hence $\zeta_A(g,h)=(\bigotimes_vT_v)\zeta_{A^0}(g,h)$,
    which transports the equality to the original physical tensors.
\end{proof}

\begin{definition}[Class-labelled four-block closures]%
    \label{def:peps_dependent_closure_class}
    \lean{TNLean.PEPS.DependentTorus.closureClass,
        TNLean.PEPS.DependentTorus.closureClass_pairConjugacyClass,
        TNLean.PEPS.DependentTorus.range_closure_commuting_eq_range_closureClass}
    \leanok
    \uses{thm:peps_dependent_closure_invariant}
    For a simultaneous conjugacy class $C=[(g,h)]$, define
    $\zeta_A(C)=\zeta_A(g,h)$. This does not depend on the representative.
    The set of closures indexed by commuting pairs equals the set indexed
    by commuting pair-conjugacy classes.
\end{definition}

\begin{theorem}[Independent classes for independent physical blocks]%
    \label{thm:peps_dependent_closure_independent}
    \lean{TNLean.PEPS.DependentTorus.linearIndependent_closureClass_canonical,
        TNLean.PEPS.DependentTorus.linearIndependent_closureClass,
        TNLean.PEPS.DependentTorus.linearIndependent_closureClass_commuting}
    \leanok
    \uses{def:peps_dependent_closure_class, def:peps_g_injective,
        def:peps_semiregular}
    Suppose each of the eight actual bond representations is semi-regular
    and every physical tensor is $G$-injective for its own incident action.
    Then the vectors $\zeta_A(C)$ are linearly independent as $C$ ranges
    over all pair-conjugacy classes. In particular, the commuting classes
    give a linearly independent family. The four physical alphabets and
    all eight virtual alphabets may differ independently.
\end{theorem}
\begin{proof}\leanok
    \uses{thm:peps_dependent_closure_extraction, thm:peps_dependent_common_inverse,
        thm:peps_full_four_cut_transport}
    First take the canonical tensors. Apply $E_{g,h}$, for a representative
    of any fixed class $C$, to a relation $\sum_{C'}\lambda_{C'}
        \zeta_{A^0}(C')=0$. The off-class coefficients vanish, and the remaining
    equation is $\lambda_C|G|^{-4}|C_G(g,h)|=0$.
    Thus every $\lambda_C=0$.
    For the actual tensors, choose their genuine local inverses with
    $L_vT_v=P_v$. Physical contraction gives
    $(\bigotimes_vL_v)\zeta_A(C)=\zeta_{A^0}(C)$ for every class $C$.
    Applying this same linear map to a relation proves the result for $A$.
    The inverse and scalar extraction are depicted below. Each vertical
    leg denotes a complete configuration, retaining its independently
    typed component indices; $\mathcal L=\bigotimes_vL_v$.
    % Formula/source: E_{a,b} (tensor_v L_v) zeta_A(g,h)
    % = E_{a,b} zeta_{A^0}(g,h), SCP10 paper_v3.tex:1590--1610,
    % the boundary-lin-indep diagram, with the eight actual trace-dual
    % factors of Lemma 4.6 applied after the four genuine local inverses.
    % Ink-to-index: bottom ket is the actual four-tensor closure zeta_A;
    % the middle box is tensor_v L_v; the top covector is E_{a,b}.
    % First contracted leg is s=(s_v)_v, s_v in Y_v. Second is the
    % complete incidence tuple sigma=(sigma_v)_v, sigma_v in X_v.
    % Right-hand leg is that same tuple sigma. Each X_v is the product
    % of its four separate incidence alphabets; no common bond type is used.
    % All legs are contracted, so the boundary signature is scalar on
    % both sides: zero virtual open indices and zero physical open indices.
    \begin{center}
        \tenkzkernel
        \begin{tenkz}[rows={bra,op,ket}, cols=1, bonds=none]
            \tn[at=(1,1), name=extractActual, size=l,
                ports={270:physical}]{E_{a,b}}
            \tn[at=(2,1), name=inverseActual, skin=mpo, size=l,
                ports={90:physical,270:physical}]{\mathcal L}
            \tn[at=(3,1), name=closureActual, size=l,
                ports={90:physical}]{\zeta_A(g,h)}
            \tnwire[name=extractedIncidences]{extractActual.270}{inverseActual.90}
            \tnmark[form=label, label pos=e]{on extractedIncidences .5}{$\sigma$}
            \tnwire[name=originalPhysical]{inverseActual.270}{closureActual.90}
            \tnmark[form=label, label pos=e]{on originalPhysical .5}{$s$}
        \end{tenkz}
        \quad$=$\quad
        \begin{tenkz}[rows={bra,ket}, cols=1, bonds=none]
            \tn[at=(1,1), name=extractCanonical, size=l,
                ports={270:physical}]{E_{a,b}}
            \tn[at=(2,1), name=closureCanonical, size=l,
                ports={90:physical}]{\zeta_{A^0}(g,h)}
            \tnwire[name=canonicalIncidences]{extractCanonical.270}{closureCanonical.90}
            \tnmark[form=label, label pos=e]{on canonicalIncidences .5}{$\sigma$}
        \end{tenkz}
        \quad$\displaystyle=|G|^{-4}
            \sum_{x\in G}\mathbf 1_{a=xgx^{-1},\,b=xhx^{-1}}$.
    \end{center}
\end{proof}

\begin{theorem}[Exact dimension of the four-cut intersection]%
    \label{thm:peps_dependent_four_cut_dimension}
    \lean{TNLean.PEPS.DependentTorus.finrank_commutingClosureSpan,
        TNLean.PEPS.DependentTorus.finrank_fourCutSpace}
    \leanok
    \uses{def:peps_full_four_cut_space, def:peps_dependent_closure_class}
    Under the hypotheses of
    Theorem~\ref{thm:peps_dependent_closure_independent}, the commuting
    class-labelled closures form a basis of $\mathcal I_A=\mathcal Z_A$.
    Both spaces have complex dimension $K$, the number of simultaneous
    conjugacy classes of commuting pairs in $G$.
\end{theorem}
\begin{proof}\leanok
    \uses{thm:peps_four_cut_closure, thm:peps_dependent_closure_independent}
    The four-cut closure theorem gives
    $\mathcal I_A=\mathcal Z_A=\spn_{\C}\{\zeta_A(C):C\text{ commuting}\}$.
    The stated family is linearly independent, so the dimension of its
    span is its cardinality $K$.
\end{proof}
